\documentclass[11pt]{article}
\usepackage[margin=1in]{geometry}
\usepackage{amsmath,amssymb}
\usepackage{booktabs}
\usepackage{graphicx}
\usepackage{hyperref}
\usepackage{xcolor}
\usepackage[numbers]{natbib}

\newcommand{\todo}[1]{\textcolor{red}{[TODO: #1]}}
\newcommand{\mdes}{MDES}

\title{Measure, Don't Search: Measurement-Driven Exit Selection for
Hardware-Constrained Early-Exit Networks Across Architectures and Devices}

\author{Airlangga Wicaksono\\ \todo{affiliation, email, co-authors}}
\date{Working draft — \today}

\begin{document}
\maketitle

\begin{abstract}
Early-exit neural networks (EENNs) trade accuracy for latency, energy, and
memory at inference time, which makes them attractive for edge deployment.
Recent work such as NACHOS designs hardware-constrained EENNs with neural
architecture search (NAS), but relies on \emph{estimated} cost models, targets
a single architecture family (CNN classifiers), and validates on a single
device class. We take the opposite approach: instead of searching over
architectures with predicted costs, we \emph{measure} the per-exit
quality/latency/energy/memory profile of already-trained multi-exit models and
reduce hardware-constrained EENN design to a cheap post-training
\emph{selection} problem. We present (i) the first cross-architecture,
cross-device empirical characterization of early-exit Pareto fronts, covering
four model families --- an encoder LM (ElasticBERT-large, 24 exits, GLUE), a
vision transformer (ViT-large, 24 exits, CIFAR-10/100), an object detector
(YOLOv9 GELAN-m, 6 exits $\times$ 3 scales = 18 anytime points, COCO), and a
decoder LLM (Llama-3.2-1B, 16 exits, generation and MCQ suites) --- trained
under two uniform self-distillation protocols (star-topology pairwise
distillation and chain-topology SEGD) and profiled on both a datacenter GPU
(NVIDIA A100) and an edge SoC (NVIDIA Jetson Orin Nano), with hardware energy
counters and power-rail telemetry rather than analytic proxies; and (ii)
\mdes{}, a measurement-driven exit-selection method that picks an exit subset
and confidence thresholds to maximize task quality under user latency, energy,
or memory budgets, solved exactly in seconds on the measured tables. Our
preliminary data shows that per-exit cost rankings are \emph{not} stable
across devices (unified-memory contention and DVFS on the edge SoC reorder
exits that are equivalent on the A100), implying that selections tuned on
datacenter hardware transfer imperfectly and motivating on-device profiling.
\mdes{} recovers hardware-fitting configurations at a
$\sim$$10^{4}\times$ lower design cost than NAS-based approaches, while
supporting architectures (detectors, decoder LLMs) that current EENN-NAS
methods do not handle.
\end{abstract}

\section{Motivation and gap}

EENNs attach auxiliary classifiers (\emph{exits}) to intermediate layers so
inference can stop early for easy inputs, or be truncated under a hard
resource budget (\emph{anytime} inference). The design-for-hardware literature
has converged on a specific paradigm --- \emph{search or co-design with
estimated costs, on CNN classifiers} --- whose blind spots our existing
infrastructure is unusually well placed to expose.

\subsection{Landscape (2020--2026)}

\begin{table}[h]
\centering\scriptsize
\begin{tabular}{lllll}
\toprule
Work & Approach & Families & Cost source & Devices\\
\midrule
HAPI'20 \citep{hapi2020} & exit placement search & CNN cls & measured latency & 1 (mobile/GPU)\\
NACHOS'24 \citep{nachos2024} & NAS & CNN cls & differentiable \emph{estimates} & 1 (est.)\\
Zniber+'25 \citep{zniber2025} & HW--algo co-design & CNN cls & \emph{analytical} DSE & simulated accel.\\
AEBNAS'25 \citep{aebnas2025} & HW-aware NAS (branches) & CNN cls & estimates & edge (single)\\
E4'25 \citep{e42025} & EE + DVFS control & CNN video & measured power & 1 edge\\
CalexNet'25 \citep{calexnet2025} & branch training + calib. & CNN cls & measured GPU lat. & 1 GPU\\
EERO'24 \citep{eero2024} & threshold/budget control & cls & FLOPs proxy & ---\\
LayerSkip'24 \citep{layerskip2024} & EE decoding & LLM & --- & datacenter\\
Diminishing'26 \citep{diminishing2026} & EE-decoding critique & LLM & --- & datacenter\\
\textbf{Ours} & measure + select & \textbf{enc-LM, ViT, detector, dec-LLM} & \textbf{measured (energy ctr.)} & \textbf{A100 + Orin Nano}\\
\bottomrule
\end{tabular}
\caption{Hardware-aware early-exit work. No prior study crosses architecture
families under one protocol, and no design method uses measured cross-device
cost tables.}
\end{table}

\subsection{Three confirmed gaps}

\begin{enumerate}
\item \textbf{Estimated-cost monoculture.} Every hardware-constrained EENN
\emph{design} method --- NACHOS \citep{nachos2024}, the multi-core co-design of
Zniber et al.\ \citep{zniber2025} (explicitly ``analytical design space
exploration''), AEBNAS \citep{aebnas2025} --- feeds the optimizer estimated or
analytical costs. Zniber et al.\ themselves show that \emph{small} architectural
variations cause \emph{disproportionate} hardware-efficiency changes via tensor
alignment and dataflow effects --- which is precisely the regime where analytic
models err. Nobody has quantified, on commodity hardware, how far
estimate-based per-exit cost rankings deviate from measured ones, nor whether
rankings are even stable \emph{across} devices. Our two-device measured tables
answer this directly, and our transfer-regret experiment (optimize on A100,
deploy on Jetson) puts a number on the cost of designing on the wrong device.
\item \textbf{CNN-classifier tunnel vision.} Both surveys
\citep{laskaridis2021adaptive, csur2024survey} state that early-exit research
concentrates on CNN image classification and Transformer NLP, with detection
and other tasks ``largely unexplored''; every entry in the table above except
the LLM-decoding line is an image classifier. Per-family methods exist
(ElasticBERT \citep{elasticbert2022}, MSDNet \citep{msdnet2018}, AnytimeYOLO
\citep{anytimeyolo2025}, LayerSkip \citep{layerskip2024}), but no work trains
four families under one protocol and profiles them under one harness --- so
the field has no answer to ``do detector exits behave like classifier
exits?'' We do (they do not: FPN-stage cliffs vs.\ smooth per-layer decay).
\item \textbf{Selection is unclaimed territory.} Threshold control exists
(EERO \citep{eero2024}, performance-control methods
\citep{performancecontrol2024}); exit-branch pruning exists
\citep{prunee2022}; serving-side batching exists \citep{drex2025}. But
\emph{joint exit-subset + threshold selection over measured multi-resource
cost tables, per device}, with offline confidence-replay evaluation, is not
published. It is the natural ``deployment-time'' complement to design-time
NAS --- and because our harness logs per-sample rows, we can evaluate any
configuration without re-running inference.
\end{enumerate}

\noindent\textbf{Positioning w.r.t.\ the LLM-skepticism result.} The
``diminishing returns'' critique \citep{diminishing2026} shows early-exit
\emph{decoding} helps less on modern, deeper-trained LLMs. Our scope differs:
edge-scale 1B models (where exits target energy/latency budgets, not just
tokens/s), task-suite quality (MCQ vs.\ generation) rather than decoding
quality alone, and measured energy. If our data confirms their pessimism for
generation while showing MCQ robustness, that is a publishable finding, not a
contradiction.

\noindent\textbf{Thesis.} When a multi-exit model already exists (trained once,
cheaply, on a frozen pretrained backbone), hardware-constrained EENN design
degenerates from architecture search to \emph{subset selection over measured
per-exit cost tables} --- solvable exactly, per device, in seconds, with
measured rather than estimated costs. The interesting science is then in the
measurements themselves: how Pareto fronts differ across architecture families,
devices, and training protocols.

\section{What already exists (assets)}

All of the following is implemented and producing data in our monorepo:

\begin{itemize}
\item \textbf{Four multi-exit model families, one training framework.} A
shared self-distillation protocol (\texttt{plan.py}: supervise-deepest, then
distill shallower exits) instantiated for ElasticBERT-large (24 per-layer
exits, LoRA), ViT-large (24 exits, LoRA), YOLOv9 GELAN-m (6 multi-scale
DDetect exits seeded from the pretrained detection head, full-head training,
frozen backbone), and Llama-3.2-1B (16 exits, LoRA + lm-head projections).
Two protocol topologies: \emph{pairwise} (star: every exit distills from the
deepest) and \emph{SEGD} (chain: exit $k$ distills from exit $k{+}1$)
\citep{loraexit2024}.
\item \textbf{Uniform profiling harness.} Per-sample CUDA-synced latency;
energy from the NVML on-die counter (A100) or from jtop/INA3221 power rails
(Jetson), with an explicit validity guard against silent telemetry loss;
peak VRAM and board-wide unified memory; SM/EMC clocks; temperature; ECE
calibration per exit. Identical JSON schema across backends and devices feeds
one exporter (CSV/plots/xlsx), including YOLO sub-exit granularity.
\item \textbf{Two devices.} NVIDIA A100-40GB (datacenter) and NVIDIA Jetson
Orin Nano 8\,GB (edge, unified memory), with an edge-hardened runner
(per-backend process isolation, OOM forensics, resume-safe sweeps).
\item \textbf{Quality metrics are the official ones per task}: GLUE metrics
via HuggingFace, top-1/top-5, COCO mAP via the reference \texttt{ap\_per\_class},
ROUGE-L / GSM8K exact-match / lm-eval-style \texttt{acc\_norm}.
\end{itemize}

\section{Research questions}

\begin{description}
\item[RQ1 (cross-architecture).] How do quality--latency--energy Pareto fronts
of early exits differ across the four families? Hypothesis: encoder LMs and
ViTs degrade smoothly per layer; detectors degrade in large steps tied to FPN
stages (scale-level sub-exits smooth this); decoder LLMs show
task-dependent cliffs (MCQ robust, generation fragile).
\item[RQ2 (cross-device, cross-power-mode).] Are per-exit cost \emph{rankings}
preserved (a) between A100 and Orin Nano, and (b) on the \emph{same} Orin Nano
across nvpmodel power caps (7\,W / 15\,W / MAXN)? Power modes rescale CPU, GPU
(SM), and memory (EMC) clocks \emph{non-uniformly}, so memory-bound and
compute-bound exits should shift relative cost --- meaning a single per-device
cost table may already be insufficient. E4 \citep{e42025} uses DVFS as a
\emph{control knob} for one CNN; no work characterizes how power modes reorder
the exit Pareto set. Corollary: how wrong is a FLOPs-only or A100-only cost
model when deployed on the edge device at a given power cap?
\item[RQ3 (training protocol).] Do pairwise and SEGD produce different
\emph{Pareto-optimal exit sets} (not just different average quality)? I.e.,
does protocol choice change which exits one should deploy?
\item[RQ4 (selection).] Given measured tables, can exact selection (\mdes{})
match or beat NAS-designed EENNs' hardware fit at negligible design cost, and
does a selection computed from A100 measurements transfer to the Jetson, or is
on-device profiling necessary?
\end{description}

\section{\mdes{}: measurement-driven exit selection}

\textbf{Inputs.} For device $d$ and exit $e \in \{1..N\}$ (sub-exits are
distinct items for the detector): measured cost vector
$c_d(e) = (\ell_d(e), j_d(e), m_d(e))$ (latency, energy/sample, peak unified
memory) and quality $q(e)$; optionally per-exit confidence distributions for
dynamic exiting.

\textbf{Static (anytime) variant.} Choose exit set $S$, $|S| \le K$, and a
deadline map; maximize terminal quality subject to
$\ell_d(\max S) \le B_\ell$, $j_d \le B_j$, $m_d \le B_m$. With measured
monotone costs this is a small exact knapsack / DP over $N \le 24$ items ---
milliseconds to solve, per budget, per device.

\textbf{Dynamic variant.} Additionally choose per-exit confidence thresholds
$\tau_e$ (entropy or max-prob). Expected cost and quality follow directly from
the logged per-sample confidence traces (no re-inference needed): we replay
recorded confidences to evaluate any $(S, \tau)$ offline, then grid/DP-optimize.
This reuse of per-sample logs is only possible because the harness stores
per-sample rows, not just aggregates.

\textbf{Transfer experiment (the headline).} Optimize $(S,\tau)$ on A100
tables, deploy on Jetson, measure regret vs.\ Jetson-optimized $(S,\tau)$.
Repeat across nvpmodel power modes (7\,W/15\,W/MAXN) to quantify how much
on-device, on-power-mode profiling buys.

\section{Preliminary evidence (to be filled from existing logs)}

\begin{table}[h]
\centering\small
\begin{tabular}{lllll}
\toprule
Family & Exits & Quality span (first$\to$last exit) & A100 lat.\ span & Jetson lat.\ span\\
\midrule
ElasticBERT-large & 24 & \todo{fill from results/bert} & \todo{} & \todo{}\\
ViT-large & 24 & \todo{fill from results/vision} & \todo{} & \todo{}\\
YOLOv9 GELAN-m & 6$\times$3 & \todo{fill from results/yolo} & \todo{} & \todo{}\\
Llama-3.2-1B & 16 & \todo{fill from results/llama} & \todo{} & \todo{}\\
\bottomrule
\end{tabular}
\caption{Scope of the measured grid. Each cell backed by
\texttt{logs/benchmark/\{model\}/...} on both devices; energy and memory
columns analogous. \todo{insert numbers via compare\_devices.py xlsx}}
\end{table}

Anecdotal observations already visible in the raw logs (to verify and
quantify): (a) Jetson unified-memory pressure penalizes deep exits
super-linearly relative to A100; (b) YOLO sub-exits at P5-only are
latency-cheap but quality-collapse on small objects, giving the detector a
much more jagged front than classifiers; (c) telemetry loss (jtop) silently
zeroes energy if unguarded --- our validity-tagging turned this from a data
corruption source into a re-run trigger. \todo{verify each on full sweep}

\section{Contributions (target claims)}

\begin{enumerate}
\item First cross-architecture (encoder LM / ViT / detector / decoder LLM),
cross-device (datacenter GPU / edge SoC) measured characterization of
early-exit quality--latency--energy--memory Pareto fronts, with per-exit
calibration; released as open data + harness.
\item Evidence that per-exit cost rankings invert between devices and power
modes, quantifying the error of FLOPs- or estimate-based EENN cost models
(direct critique of the cost-model assumption in NAS-based design).
\item \mdes{}: an exact, seconds-scale, measurement-driven exit + threshold
selection method under multi-resource budgets, including an offline
confidence-replay evaluator; comparison of design cost and achieved fit
vs.\ NACHOS-style NAS.
\item A protocol study (pairwise vs.\ SEGD) showing whether training topology
changes the deployable Pareto set.
\end{enumerate}

\section{What's next (execution plan)}

\begin{enumerate}
\item \textbf{Finish the measured grid} (blocking): retrain YOLO with the
fixed full-head protocol; complete Jetson sweeps for all four backends
(\texttt{bench\_jetson.py all -d}); run \texttt{verify --real} until all
leaves pass.
\item \textbf{Power-mode ablation}: repeat the Jetson HW sweep under
nvpmodel 7W/15W/MAXN; the harness already logs nvpmodel in device caps.
\item \textbf{Implement \mdes{}} ($\sim$300 lines over existing JSONs): cost
tables $\to$ DP selection $\to$ confidence-replay for dynamic thresholds;
output = chosen exit sets per budget per device.
\item \textbf{Transfer experiment}: A100-optimized vs Jetson-optimized
selections, regret table; the paper's central figure.
\item \textbf{Baselines}: (a) FLOPs-proportional cost model (what NAS
methods assume) fed to the same selector; (b) uniform exit spacing;
(c) single-exit (no EE) per budget; (d) reported NACHOS numbers
(CIFAR/SVHN scope) for design-cost comparison.
\item \textbf{Writing}: fill Tables/Figures from \texttt{compare\_devices.py}
outputs; related-work pass; submit.
\end{enumerate}

\section{Venue and positioning (min Q3, expected Q2)}

\begin{itemize}
\item \textbf{Primary}: Journal of Systems Architecture (Elsevier, Q1/Q2
systems) --- fits measurement + edge deployment + method.
\item \textbf{Alternatives}: ACM TECS (Q2, embedded), Future Generation
Computer Systems (Q1, if we emphasize the edge-cloud continuum),
Sustainable Computing: Informatics and Systems (Q2, energy angle),
IEEE Internet of Things Journal (Q1, stretch).
\item \textbf{Fallback (Q3)}: Microprocessors and Microsystems, or
IEEE Access (Q2 but high-volume).
\end{itemize}
Positioning against NACHOS: complementary, not competing --- NAS designs the
network; we show that once exits exist, \emph{deployment-time} constraint
satisfaction is a measurement + selection problem, and that measurements are
non-transferable enough across devices that search-time cost models are
structurally optimistic.

\bibliographystyle{unsrtnat}
\bibliography{references}

\end{document}
